%/*% pdflatex -file-line-error -interaction nonstopmode %

\documentclass{article}
\usepackage{graphicx}
\usepackage{amsmath}
\usepackage{amsfonts}
\usepackage{bookmark}
\usepackage{booktabs}
\usepackage{exercises}
\graphicspath{ {./images/test} }
\usepackage{indentfirst}
\usepackage{amsthm}
\usepackage{yhmath}
\usepackage{placeins}
\newcommand*{\oldneg}{\mathord{\sim}}
\newcommand\showdiv[1]{\overline{\smash{)}#1}}
\newcommand\Mydiv[1]{%
\smash{\raise.3ex\hbox{$\big)$}}\mkern-8mu
        \overline{\enspace\strut#1}}
\setlength{\tabcolsep}{5pt}
\theoremstyle{definition}
\newtheorem{definition}{Definition}

\setlength{\parindent}{.5in}
\setlength{\parskip}{12pt}
\renewcommand{\arraystretch}{2}

\title{3.3 Answers}
\begin{document}
\maketitle

\section*{1}
The function $h(x)=g(f(x))$ is called the composition of the function $g$ with the function $f$.

\section*{2}
The sum of two power functions is $sometimes$ a power function. For example, the sum of $ax^n$ and $bx^n$ is $ax^n+bx^n=(a+b)x^n$, which is a power 
function, whereas the sum of $x^2$ and $2x^3$ is $2x^3+x^2$, which is not a power funciton.

\section*{8}
\[ g(x)-f(x)=(5x^\frac{1}{3}+4x^\frac{1}{2})-(-3x^\frac{1}{3}+4x^\frac{1}{2}) \]
\[ =8x^\frac{1}{3} \]
Because a cube root is defined for all real numbers, the domain of $g(x)-f(x)$ is all real numbers. However, because both $f(x)$ and $g(x)$ involve 
taking square roots, the domain is actually all real non-negative numbers.

\section*{10}
\[ g(x)-g(x)=(5x^\frac{1}{3}+4x^\frac{1}{2})-(5x^\frac{1}{3}+4x^\frac{1}{2})=0 \]
Because $0$ does not rely on the value of $x$, the domain of $g(x)-g(x)$ is any real number. However, because $g(x)$ involves 
taking square roots, the domain is actually all real non-negative numbers.

\section*{11}
\[ f(x)+g(x)=-7x^\frac{2}{3}-1+2x^\frac{2}{3}+6=-5x^\frac{2}{3}+5 \]
B

\section*{13}
\[ g(x)\times f(x)=5x^\frac{1}{2}\times 4x^\frac{2}{3}=20x^{\frac{1}{2}+\frac{2}{3}}=20x^\frac{7}{6} \]
The domain is all non-negative real numbers, as negative real numbers do not have a square root.

\section*{17}
\[ \frac{g(x)}{f(x)}=\frac{5x^\frac{1}{2}}{4x^\frac{2}{3}}=\frac{5x^{\frac{1}{2}-\frac{2}{3}}}{4}=\frac{5x^{-\frac{1}{6}}}{4}=\frac{5}{4x^\frac{1}{6}} \]
\[ =\frac{5x^\frac{5}{6}}{4x} \]
The domain is all positive real numbers, as negative real numbers do not have a square root, and we cannot divide by $4(0)=0$.

\section*{22}
\[ h(f(-9))=h(3(-9)+2)=h(-25)=\frac{-25-2}{5}=-\frac{27}{5} \]

\section*{23}
\[ g(h(8))=g(\frac{8-2}{5})=g(\frac{6}{5})=-(\frac{6}{5})^2=-\frac{36}{25} \]

\section*{26}
\[ h(h(-4))=h(\frac{-4-2}{5})=h(-\frac{6}{5})=\frac{-\frac{6}{5}-2}{5}=-\frac{\frac{16}{5}}{5}=-\frac{16}{25} \]

\section*{28}
\[ g(f(x))=g(3x^{-1})=2(3x^{-1})-7=6x^{-1}-7 \]
The domain is all real numbers other than $0$, as $0^{-1}$ is undefined.

\section*{31}
\[ g(h(x))=g(\frac{x+4}{3})=2(\frac{x+4}{3})-7=\frac{2x+8}{3}-7 \]
The domain is all real numbers, as the domains of both $g$ and $h$ are all real numbers.

\section*{37}
The error here is that $g(x)$ was mistaken to be the function $4x-3$. Instead, it is the function $4x$, meaning 
\[ g(f(x))=g(x^2-3)=4(x^2-3)=4x^2-12 \]

\section*{38}
\[ g(f(x))=g(7x^2)=3(7x^2)^{-2}=3(49^{-1}x^{-4})=\frac{3}{49x^4} \]
A

\section*{39}
We can use the commutative property of multiplication to determine that, if $f=ax$ and $g=bx$, then $f(g(x))=abx$ and $g(f(x))=abx$, 
meaning $f(g(x))=g(f(x))$.

\section*{42}
$f(x)=|x|,g(x)=2x+9$

\section*{44}
$C(x(t))=C(50t)=60(50t)+750=3000t+750$ and \[ C(x(5))=C(50(5))=C(250)=60(250)+750=15750 \]
This number means that, in $5$ hours, the company produces $250$ sneakers, which costs $\$15750$.

\section*{45}
(a): We know that the functions for taking the $\$15$ discount and the $10\%$ discount are $f(x)=x-15$ and $g(x)=x-0.1x=0.9x$, respectively. To find the 
sale price when the $\$15$ discount is taken first, we find:
\[ g(f(x))=g(85-15)=g(70)=0.9(70)=63 \]
(b): Similarly, to find the sale price when the $10\%$ discount is taken first, we find:
\[ f(g(x))=f(0.9(85))=f(76.5)=76.5-15=61.5 \]
(c): The order of taking the $10\%$ discount first gives the better deal, as $61.5<63$.

\section*{46}
(a): We know that the dog is swimming at $0.9$ mmeters per second. We also know, by the Pythagorean Theorem, that the distance he is swimming is 
$\sqrt{x^2+144}$, meaning a function representing the time the dog is swimming is given by $s(x)=\frac{\sqrt{x^2+144}}{0.9}$ seconds. 
We know that the dog ran $20-x$ meters in 6.4 meters per second, meaning this took the dog $r(x)=\frac{20-x}{6.4}$ seconds.\newline
(b): The total time spent traveling is given by the sum of the above two functions, that is:
\[ t(x)=\frac{\sqrt{x^2+144}}{0.9}+\frac{20-x}{6.4} \]
(c): The value of $x$ that minimizes $t(x)$ is $1.704$, yielding a value of $16.326$. This means that, when the value of $x=1.704$, the dog travels for a 
combined time of $16.326$ seconds.

\end{document}
